\documentclass[11pt]{article}
\usepackage[margin=1.1in]{geometry}
\usepackage{amsmath,amssymb,amsthm}
\usepackage[hidelinks]{hyperref}
\newtheorem{theorem}{Theorem}
\title{The Feynman--Kac Formula via Random Walks}
\author{}
\date{}
\begin{document}
\maketitle

\begin{abstract}
A nearest-neighbor walk gives a discrete path formula for a heat equation
with potential. The functional central limit theorem takes the formula to
Brownian motion, and the Markov property identifies the limit with the
unique bounded mild solution.
\end{abstract}

\begin{theorem}
Let $f,q$ be bounded continuous real functions on $\mathbb R^d$ and set
$\Delta=-\sum_{j=1}^d\partial_j^2$. The unique mild solution, bounded on
each finite time strip, of
\[
  \partial_tu+\Delta u+qu=0,\qquad u(0,x)=f(x)
\]
is given by
\[
  u(t,x)=\mathbb E_x\!\left[
   f(B_t)\exp\!\left(-\int_0^t q(B_s)\,ds\right)\right],
  \qquad B_s=x+\sqrt{2}\,W_s,
\]
where $W$ is standard $d$-dimensional Brownian motion.
\end{theorem}

\section{The discrete path formula}
For a time step $\tau>0$, put $h=\sqrt{2d\tau}$. Start a walk at
$X_0=x$, and let each step be independently uniform on
$\{\pm h e_j:1\le j\le d\}$. Denote its one-step averaging operator
by $A_\tau$ and set $\Delta^\tau=(I-A_\tau)/\tau$.
Taylor expansion on smooth test functions gives
$\Delta^\tau\varphi(x)\to\Delta\varphi(x)$. The covariance of one
increment is $(h^2/d)I=2\tau I$, which fixes the Brownian clock.

Evaluate the potential term at the new time level and the discrete
Laplacian at the old one. For $\tau\|q\|_\infty\le 1/2$ the scheme gives
\[
  \frac{u_{m+1}-u_m}{\tau}+\Delta^\tau u_m+q\,u_{m+1}=0,
  \qquad
  u_{m+1}(x)=\frac{(A_\tau u_m)(x)}{1+\tau q(x)}.
\]
Starting from $u_0=f$, induction by conditioning on $X_1$ yields
\[
  u_m(x)=\mathbb E_x\!\left[
     f(X_m)\prod_{k=0}^{m-1}(1+\tau q(X_k))^{-1}\right].
\]
The time step contributes a factor at the departure point.
There are $m$ factors for $m$ steps. For $f\equiv1$ and $q\equiv c>0$,
the scheme gives $(1+\tau c)^{-m}$, which tends to $e^{-ct}$ at time $t$.

\section{The Brownian limit}
Fix $t>0$ and set $\tau=t/n$. Let $Y^\tau$ be the piecewise-constant
walk, equal to $X_k$ on $[k\tau,(k+1)\tau)$ and to $X_n$ at $t$.
By the functional central limit theorem the linear interpolation
converges in law in $C([0,t];\mathbb R^d)$ to $B$. The uniform
distance from $Y^\tau$ to that interpolation is at most $h$.
Equivalently, $Y^\tau$ converges to $B$ in the Skorokhod topology,
where convergence to a continuous path implies uniform convergence.

Since $q$ is bounded, the logarithm has a uniform quadratic
remainder. Thus, for every walk path,
\[
 -\sum_{k=0}^{n-1}\log(1+\tau q(X_k))
 =-\int_0^t q(Y^\tau_s)\,ds+O(t\tau).
\]
At each continuous path, endpoint evaluation and the time integral of
$q$ are continuous under uniform convergence. For the integral, uniform
convergence confines the paths to a common compact neighborhood of the
limiting path, on which $q$ is uniformly continuous. The continuous
mapping theorem therefore applies to the endpoint and the integral
jointly. The weighted endpoint values are bounded in absolute value by
$\|f\|_\infty e^{2t\|q\|_\infty}$ because
$|\log(1+\tau q)|\le 2\tau\|q\|_\infty$. Their expectations converge to
the expectation in the theorem.

\section{Identification with the heat equation}
Write $P_t=e^{-t\Delta}$ for the free heat semigroup and let $u$ denote
the limiting expectation. Pathwise differentiation with respect to
the lower endpoint of the integral gives
\[
 e^{-\int_0^t q(B_r)\,dr}
  =1-\int_0^t q(B_s)e^{-\int_s^t q(B_r)\,dr}\,ds.
\]
Multiply by $f(B_t)$, take expectations, and condition on $B_s$ inside
the integral. Boundedness justifies Fubini, and the Markov property
gives
\[
 u(t)=P_tf-\int_0^t P_s\bigl(q\,u(t-s)\bigr)\,ds
     =P_tf-\int_0^t P_{t-s}\bigl(q\,u(s)\bigr)\,ds.
\]
This is the mild equation. For two bounded solutions on a finite time
strip, the sup norm of their difference at time $t$ is bounded by
$\|q\|_\infty$ times the integral of that norm from $0$ to $t$.
Gr\"onwall's inequality gives uniqueness. With smooth bounded data
and bounded derivatives, the mild solution is classical by standard
heat equation regularity. \hfill$\square$

\begin{thebibliography}{9}
\bibitem{taylor} M.~E.~Taylor, \emph{Brownian Motion and Potential
Theory}, Chapter 11, Sec.~2.
\href{https://mtaylor.web.unc.edu/wp-content/uploads/sites/16915/2018/04/chap11.pdf}{Online chapter}.
\end{thebibliography}
\end{document}
